\section{Approach \& Uniqueness}

\subsection{Getting Started}

Since much of the Embedded Xinu kernel relied on the specific hardware state that the boot loader created before its crash, we were determined to refactor into the proper privilege mode with as few changes as possible, therefore limiting the damage that restructured hardware permissions could cause. This refactor required two main efforts; firstly we needed to rebuild the boot loader such that it would complete its task and safely hand off to the kernel; secondly we needed to modify the kernel so that it would properly handle its new privilege level and permissions, ideally without disrupting more components than those in the initialization process. 

\subsection{The Boot Loader}

The boot loader used on the Embedded Xinu system had two major flaws. The first flaw, as previously mentioned, was that the boot loader would crash somewhere in the U-boot step, and let the kernel have access to machine mode. The second flaw was that the boot loader was bundled with a 10 gigabyte Linux distribution, which acted as a fallback for the Xinu kernel. In order to address these issues we had to find an updated version of U-boot for RISC-V, and build it from source. While reading into the U-boot documentation, we discovered that much of the hardware configuration done by our network booting script could in fact be hard-coded into the boot loader configuration. This along with disabling scanning for the unused USB devices allowed us to not only correct the boot sequence, but shave off nearly 20 seconds from every run.

\subsection{Initialization \& Interrupt Refactoring}

Once the boot loader was rebuilt, we ran into issues with both the kernel initialization sequence and with interrupt handling. Since the kernel now started in S-mode (as opposed to M-mode), all the initialization around the M-mode control and status registers (CSRs) were disallowed by the hardware. To resolve this we had to root out any assembly calls that modify M-mode privileged registers, and ensure that all the configuration changes originally made by the kernel are made by the rebuilt boot loader. With the kernel now able to boot past the initialization phase, we had to sort out exception trapping. While the existing exception handling infrastructure was viable, the privilege change prevented them from being trapped by the kernel handler. This was resolved with CSR reconfiguration and a deeper look at the RISC-V interrupt and exception structure. Finally we implemented timer control through the platform level interrupt controller (PLIC). The PLIC allows for hardware outside the CPU to raise interrupts that can be trapped and handled like standard kernel level interrupts. In this case we configured the external timer module to raise an interrupt every second, thus allowing the kernel to track time at that interval.
